\clearpage
\section{Nazwa i opis projektu}		%1

Projekt nosi nazwę \textbf{System wykrywania intruzów}. Jego celem jest zwiększenie bezpieczeństwa obiektu poprzez automatyczną analizę obrazu z kamer monitoringu, plików wideo oraz strumieni sieciowych. System wspomaga operatora ochrony przez wykrywanie osób, klasyfikowanie ich jako pracowników lub potencjalnych intruzów oraz zapisywanie materiału dowodowego w przypadku wystąpienia zdarzenia alarmowego.

System działa w oparciu o źródła obrazu rozmieszczone w wybranych strefach obiektu, takich jak wejścia, korytarze, magazyny, parkingi oraz strefy ograniczonego dostępu. Obraz jest analizowany na bieżąco, a wykryte zdarzenia są klasyfikowane zgodnie z trybem pracy, aktualnymi progami alarmu oraz konfiguracją modelu.

Istotnym elementem projektu jest rozróżnienie sposobu działania systemu w zależności od pory dnia. W trybie dziennym aplikacja wykorzystuje model segmentacyjny YOLO, który wyznacza maskę sylwetki wykrytej osoby. Na podstawie tej maski pobierane są kolory górnej oraz dolnej części ubioru, a następnie porównywane z zaprogramowanym wzorcem stroju pracownika. Takie podejście ogranicza wpływ tła na analizę koloru i pozwala odróżnić pracowników od osób niepożądanych. W trybie nocnym system działa bardziej restrykcyjnie: wykorzystuje klasyczny model detekcji YOLO i traktuje każdą wykrytą osobę jako potencjalne naruszenie bezpieczeństwa.

\begin{itemize}
\item \textbf{Tryb dzienny}: system analizuje maskę sylwetki osoby i porównuje kolor ubioru z ustawionym wzorcem pracownika.
\item \textbf{Tryb nocny}: każda wykryta osoba traktowana jest jako potencjalne naruszenie bezpieczeństwa.
\item \textbf{Reakcja systemu}: po wykryciu incydentu zapisywany jest czas zdarzenia, źródło obrazu, tryb pracy, liczba osób oraz materiał dowodowy w postaci klipu wideo.
\item \textbf{Obsługa przez operatora}: konfiguracja kamer, modeli, progów detekcji i archiwizacji odbywa się z poziomu graficznego interfejsu aplikacji.
\end{itemize}

Projekt stanowi bazę do dalszego rozwoju, między innymi o integrację z systemami kontroli dostępu, powiadomienia zewnętrzne, dokładniejszą analizę zachowań osób oraz rozbudowane raportowanie incydentów.
